\section{Contesto e riferimenti teorici}
\label{sec:contesto}

Questa sezione richiama, in forma sintetica, i concetti del corso di Internet
of Things che hanno guidato le scelte progettuali di PomodorIA, evidenziando
esplicitamente il collegamento fra teoria e implementazione.

\subsection{Il modello IoT World Forum (IoTWF)}
\label{sec:iotwf}

Il modello di riferimento adottato per strutturare l'architettura del sistema
è il modello a sette livelli \textbf{IoT World Forum}, che descrive un
sistema IoT come una pipeline che va dal dispositivo fisico fino ai processi
di collaborazione con l'utente finale:

\begin{enumerate}
    \item \textbf{Physical Devices \& Controllers} --- i dispositivi fisici
    (sensori, attuatori) che percepiscono e agiscono sul mondo reale.
    \item \textbf{Connectivity} --- l'infrastruttura di comunicazione tra i
    dispositivi e il resto del sistema.
    \item \textbf{Edge (Fog) Computing} --- l'elaborazione dei dati vicino
    alla fonte, per ridurre latenza e dipendenza dal cloud.
    \item \textbf{Data Accumulation} --- la persistenza dei dati raccolti.
    \item \textbf{Data Abstraction} --- la normalizzazione e strutturazione
    dei dati per renderli utilizzabili dal livello applicativo.
    \item \textbf{Application} --- la logica applicativa e le interfacce che
    rendono il sistema utilizzabile.
    \item \textbf{Collaboration \& Processes} --- l'integrazione con i
    processi umani e organizzativi (notifiche, decisioni condivise).
\end{enumerate}

Questo modello è utilizzato in maniera sistematica nella
Sezione~\ref{sec:architettura} per motivare esplicitamente ogni scelta
implementativa del progetto.

\subsection{Smart Agriculture e Agriculture 4.0}

Il corso inquadra la Smart Agriculture come uno degli ambiti applicativi
principali dell'IoT, in stretta relazione con il concetto di
\emph{Climate-Smart Agriculture}: l'uso di sensori (stazioni meteo, sensori
di umidità del suolo, videocamere, droni) per raccogliere dati su ogni
singola area o coltura, permettendo agli agricoltori di decidere in anticipo
e con precisione come distribuire acqua, pesticidi o fertilizzanti. Questa
evoluzione viene descritta come un percorso a quattro fasi:

\begin{itemize}
    \item \textbf{Agriculture 1.0} (1784--1870 circa): agricoltura basata su
    risorse umane e animali, scarsa efficienza.
    \item \textbf{Agriculture 2.0} (XX secolo): agricoltura meccanizzata, uso
    ancora inefficiente delle risorse.
    \item \textbf{Agriculture 3.0} (1992--2017): agricoltura automatizzata ad
    alta velocità, ma con scarsa intelligenza nel processo decisionale.
    \item \textbf{Agriculture 4.0} (dal 2017 ad oggi): Smart Agriculture,
    basata sull'uso di intelligenza artificiale e sistemi informativi
    moderni per analizzare, elaborare e decidere.
\end{itemize}

PomodorIA si colloca esplicitamente in quest'ultima fase: l'elemento
distintivo del progetto non è la sola raccolta di dati ambientali (già
presente nelle fasi precedenti), ma l'uso di un modello di intelligenza
artificiale (la CNN) per l'analisi diretta dello stato di salute della
pianta a partire da un dato visivo, integrato con un agente decisionale
automatico.

\subsection{Agenti razionali e modello PEAS}
\label{sec:peas}

Il corso introduce il concetto di \emph{agente razionale} come entità che
percepisce l'ambiente tramite sensori e agisce su di esso tramite attuatori,
scegliendo l'azione che massimizza una misura di prestazione attesa. La
caratterizzazione formale di un agente avviene tramite il modello
\textbf{PEAS}:

\begin{itemize}
    \item \textbf{P -- Performance}: il criterio con cui si misura il
    successo dell'agente.
    \item \textbf{E -- Environment}: l'ambiente in cui l'agente opera, e le
    sue proprietà (osservabilità, determinismo, staticità/dinamicità,
    discretezza/continuità).
    \item \textbf{A -- Actuators}: gli strumenti con cui l'agente agisce
    sull'ambiente.
    \item \textbf{S -- Sensors}: gli strumenti con cui l'agente percepisce
    l'ambiente.
\end{itemize}

Il PEAS applicato all'agente decisionale di PomodorIA è discusso nel
dettaglio nella Sezione~\ref{sec:agent-layer}, dove l'agente viene inoltre
classificato secondo le tipologie standard degli agenti (reattivo,
basato su modello, basato su obiettivi).

\subsection{Data Fusion}

Un ulteriore concetto cardine del progetto è la \textbf{Data Fusion}: la
combinazione di dati provenienti da fonti eterogenee (nel nostro caso,
un'immagine e una serie di letture ambientali numeriche) in un'unica
rappresentazione coerente, per prendere decisioni più robuste di quelle
ottenibili da una singola fonte isolata. Nel progetto, questo principio è
implementato correlando statisticamente i valori ambientali simulati alla
patologia rilevata nell'immagine, cosicché l'agente decisionale possa
ragionare su un quadro informativo coerente (ad esempio, una diagnosi di
malattia fungina accompagnata da un valore di umidità plausibilmente
elevato), esattamente come avverrebbe incrociando dati reali da sensori
fisici differenti.

\subsection{Edge e Fog Computing}

Il corso motiva lo spostamento della computazione verso i margini della
rete (Edge/Fog Computing) con la necessità di ridurre la latenza di
risposta e la dipendenza da connettività verso il cloud, un vincolo
particolarmente rilevante in scenari come le serre agricole, spesso
collocate in aree con connettività limitata. PomodorIA riflette questa
motivazione teorica nella scelta di eseguire l'intera inferenza della CNN
localmente (su CPU, con vincoli di risorse simulati), applicando tecniche
di compressione del modello (Sezione~\ref{sec:compressione}) proprio per
renderlo eseguibile su un dispositivo Edge a risorse limitate come un
Raspberry Pi.
